\homework{3}{Fri 10 Oct 2025 10:27}{}

\textbf{Algebra of Charges.}
Consider the Lagrangian density $\mathcal{L} = \partial_ \mu \phi_a^\dag \partial^ \mu\phi_a - V(\phi^a)$.
Suppose it is symmetric under action of Lie group, described by the Lie algebra
\[
	\left[ \hat{\lambda} _A, \hat{\lambda} _B \right] = i f_{AB}^c \hat{\lambda} _C
.\] 
Compute the invariant current and charges as given by the Noether's theorem. After that, quantize the fields by considering the canonical commutation relations
\[
	[\pi_a(x), \phi_b(y)] = - i \delta (x - y) \delta_{a b}
\] 
and compute the charge and current algebras.

\begin{proof}[Solution]
	We look at the action of transformation on $\phi$ and $\phi^\dag$.
	\[
		\begin{cases}
			\phi^a \to \phi^a + i \theta^A (\hat{\lambda} _A)^a_b \phi^b + O(\theta^2) \\
			\phi^\dag_a \to \phi^\dag_a - i \theta^A \phi^\dag _b (\hat{\lambda} _A)^b_a + O(\theta^2) \\
		\end{cases}
	.\] 
	The Noether current is given by
	\[
		J_A^\mu = \frac{\partial \mathcal{L}}{\partial (\partial_ \mu \phi^a)} i (\hat{\lambda} _A) ^a_b\phi^b + \frac{\partial \mathcal{L}}{\partial (\partial_ \mu \phi^\dag_a)} i \phi_b^\dag (\hat{\lambda} _A)^b_a
	.\] 
	Compute the derivatives,
	\[
		\frac{\partial \mathcal{L}}{\partial \left( \partial _ \mu \phi^a \right)} = \partial _ \mu \phi^\dag _a,
		\qquad
		\frac{\partial \mathcal{L}}{\partial \left( \partial _ \mu \phi^\dag _a \right) }  = \partial_ \mu \phi^a
	.\] 
	Hence
	\[
		J_A^ \mu=
		i\left( \partial _ \mu \phi^\dag \hat{\lambda} _A \phi - \phi^\dag \hat{\lambda} _A \partial _ \mu \phi \right) 
	.\] 
	\[
		Q_A = \int d^3 x J_A^0, \quad \frac{\partial }{\partial t} Q_A = 0
	.\] 

	Define the conjugate momentums
	\[
		\pi_a(x) = \frac{\partial \mathcal{L}}{\partial (\partial_0 \phi^a(x))} = \partial_0 \phi^\dag _a, \quad
		\pi^{\dag a}(x) = \partial_0 \phi^a
	.\] 
	In the canonical commutation relations, we only have nontrivial commutator between $\phi$ or $\phi^\dag$ and their respective conjugate momentum. That is, $[\pi, \pi^\dag] = [\phi, \phi^\dag] = [\pi^\dag, \phi] = [\pi, \phi^\dag] = 0$.
	Now we compute the commutation relations.
	\begin{align*}
		\left[ Q_A, \phi(y) \right] 
		&= \int d^4 x i\left[ (\pi \lambda_A \phi - \phi^\dag \lambda_A \pi^\dag), \phi(y) \right] \\
		&= \int d^4 x i [\pi(x), \phi(y)] \lambda_A \phi \\
		&= \lambda_A \phi(y)
	.\end{align*}
	Similarly, 
	\[
		[Q_A, \phi^\dag (y)] = \phi^\dag(y) \lambda_A, \quad
		[Q_A, \pi(y)] = - \pi(y) \lambda_A,\quad
		[Q_A, \pi^\dag(y)] = - \lambda_A \pi^\dag (y)
	.\] 
	To compute the commutation of currents, 
	\begin{align*}
		&[\pi(x) \lambda_A \phi(x), \pi(y) \lambda_B \phi(y)]\\
		&= \pi(x) \lambda_A [\phi(x), \pi(y)] \lambda_B \phi(y) + \pi(y) \lambda_B \lambda_A [\pi(x), \phi(y)] \phi(x) \\
		&= - i \delta(x - y) \left( \pi \lambda_B \lambda_A \phi - \pi \lambda_A \lambda_B \phi \right)  \\
		&= - i \delta(x - y) \pi [\lambda_B, \lambda_A] \phi \\
		&= - \delta(x - y) f_{A B}^C \pi \lambda_C \phi
	.\end{align*}
	In the first equality, we use the identity $[A B, C D]=A[B, C] D+[A, C] B D+C A[B, D]+C[A, D] B$. We also uses the fact that $\lambda_A$ commutes with the fields.
	Similarly,
	\begin{equation*}
		\left[ \phi^\dag(x) \lambda_A \pi^\dag(x), \phi^\dag(y) \lambda_B \pi^\dag(y) \right]
		= \delta(x - y) f_{A B}^C \phi^\dag \lambda_C \pi^\dag
	.\end{equation*}
	Thus,
	\begin{align*}
		[J_A^0, J_B^0] 
		&= i^2 \left( [\pi \lambda_A \phi,\lambda \lambda_B \phi] + [\phi^\dag \lambda_A \pi^\dag, \phi^\dag \lambda_B \pi^\dag] \right)  \\
		&= \delta(x - y) f_{A B}^C \left( \pi \lambda_C \phi - \phi^\dag \lambda_C \pi^\dag \right)  \\
		&= - i \delta(x - y) f_{A B}^C J_C^0
	.\end{align*}
	We see that the current algebra has same structure as the Lie algebra $\lambda_A$. Integrate this on both sides and we find the same for charge algebra.
	\[
		[Q_A, J_B^0] = - i f_{A B}^C J_C^0, \qquad
		[Q_A, Q_B] = - i f_{A B}^C Q_C
	.\] 
\end{proof}

\textbf{Non-relativistic Galilean Group.} Derive the commutation relations of the non-relativistic Galilean group generated by spatial roatations, spatial-temporal translations, and non-relativistic boosts of the form $x \to x + v t$. Show that this group admits a central extension that is the total mass of the system.

\begin{proof}[Solution]
	All Gelilean transformations are given in the form
	\[
		\begin{cases}
			 x \to (1 + \omega^i J_i) x + v^i t + a^i \\
			 t \to  t + \tau
		\end{cases}
	.\] 
	They are parameterized by $(\omega^i, v^i, a^i, \tau)$, which is 10 dimensions in total.

	In the classical case, we have the following commutation relations:	
	\begin{align}
\begin{aligned}
{\left[J_i, J_j\right] } & =i \epsilon_{i j k} J_k \\
{\left[J_i, P_j\right] } & =i \epsilon_{i j k} P_k \\
{\left[J_i, K_j\right] } & =i \epsilon_{i j k} K_k \\
{\left[P_i, H\right] } & =0, \quad\left[K_i, H\right]=i P_i, \\
{\left[K_i, P_j\right] } & =0, \quad\left[K_i, K_j\right]=0 .
\end{aligned}
\end{align}

In quantum mechanics, representation is given by the Hilbert space of wave functions $\psi(x, t)$ and all transformations must leave invariant the Schrodinger equation
\[
	i \partial_t \psi = - \frac{1}{2 m} \nabla ^2 \psi
.\] 
This enforces us to modify the action of $K$ :
\[
	\exp \left( - i K_i v^i \right) \psi(x,t)
	= e^{i ( m v x - \frac{1}{2} m v^2 t)} \psi(x + v t, t)
.\] 
All other generators keep the equation invariant and we don't have to modify.

Now, we can compute $U(v) U(a) U(v)^{-1} U(a)^{-1}$ :
\[
	\exp \left( - i v K \right) \exp \left( - i a P \right) \exp \left( i v K \right) \exp \left( i a P \right) = e^{i m v a}
.\] 
Expand and look at lowest-order terms, we have
\[
	[v K, a P] = i m a v \implies [K_i, P_j] =i m \delta_{i j} 
.\] 
Promote $m$ to a central element $M$, the commutation reads
\[
	[K_i, P_j] = i \delta_{i j} M
.\] 
The Galilean algebra with the above extension is called the Bargmann algebra.
\end{proof}

\textbf{Adjoint representation.} For the Lie algebra $[\hat{\lambda} _k, \hat{\lambda} _m] = i C_{k m}^l \hat{\lambda} _\lambda$, show that 
 \[
	 (\hat{\lambda} _k)^\lambda_m = i C_{k m}^l
\] 
is a representation and determine its dimension.

\begin{proof}[Solution]
	We are set to prove
	\[
		[\lambda_m, \lambda_n]^k_l
		\stackrel!= (i C_{m n} ^p \lambda_p)^k _l 
		= - C _{m n}^p C_{p l}^k
	.\] 
	Let's compute:
	\begin{align*}
		[\lambda_m, \lambda_n]^k_l
		&= (\lambda_m)^k_s (\lambda_n)^s_l - (\lambda_n)^k_s (\lambda_m)^s_l \\
		&= - C_{m s}^k C_{n l}^s + C_{n s}^k C_{m l}^s \\
		&= C_{n l}^s C_{s m}^k - C_{m l}^s C_{s n}^k \\
		&\triangleq (- C_{l m}^s C_{s n}^k - C_{m n}^s C_{s l}^k) - C_{m l}^s C_{s n}^k \\
		&= - C_{m s}^s C_{s l}^k
	\end{align*}
	where in the $\triangleq$ we used the Jacobi identity
	\[
		\left[ [a,b] , c \right]  + \text{cyc.} = 0
	,\] 
	i.e.,
	\[
		C_{n l}^s C_{s m}^k + C_{l m} ^s C_{s n}^k + C_{m n}^s C_{s l}^k = 0
	.\] 
	The dimension of the representation would be size of the Lie algebra.
\end{proof}

\textbf{Fundamental group of $SO(3)$.} Show that the fundamental group of $SO(3)$ is $\mathbb{Z}_2$. Explain why $SO(3)$ is not simply connected but $SO(2) $ is.

\begin{proof}[Solution]
	In general, any connected Lie group $G$ admits a universal cover
	\[
		\varphi: \hat{G}  \to  G
	\] 
	where $\hat{G} $ is simply connected, and $\pi_1(G) = \operatorname{ker} \varphi$. To show this, we explicitly construct the universal cover $\hat{G} $ by first considering the path group
	\[
		PG = \{\gamma: [0,1] \to G | \gamma \text{ smooth, }\gamma(0) = \text{Id}\} 
	.\] 
	Define on $PH$ multiplication and inverse pointwise:
	\[
		(\gamma \gamma')(t) = \gamma(t) \gamma'(t),
		\gamma^{-1}(t) = \gamma(t)^{-1}
	.\] 
	This makes $PH$ a simply connected Lie group. Now, define the covering map
	\[
		p: PH \to  H, \qquad p(\gamma) = \gamma(1)
	.\] 
	Then $\operatorname{ker} p$ gives all smooth loops based on $\text{Id} \in G$. Now quotient out the trivial homotopy class:
	\[
		\tilde p : P H / [0] \to  H
	.\] 
	And one gets the universal cover. From construction one sees that $\operatorname{ker} \tilde p = \pi_1(H)$.

	So why is $SO(3)$ not simply connected? Because it admits a universal covering $SU(2) \to SO(3)$ with the kernel being $\operatorname{ker} \varphi = \mathbb{Z}_2$, identifying antipodal points.

	$SO(2)$ is only the unit circle which for sure is simply connected.
\end{proof}





